% FAC-EN-U0069
% Source: fac_body.tex lines 2624--2696.
% Authority: PDF physical pp. 69--70; printed pp. 264--265
% (Chapter III no. 71 complete).
% French control: T0042 repairs three inherited epsilon commands to the
% printed membership sign. No printed-source correction is admitted here.

\phantomsection\label{fac:III:4:no:71}
\textbf{71. Determination of $T^r(M)$.}
(In this number, as in the following one, we shall assume that $r\geq1$;
the case $r=0$ leads to somewhat different, though trivial, statements.)

\phantomsection\label{fac:III:4:no:71:def:Omega}
We shall denote by $\Omega$ the graded $S$-module $S(-r-1)$; it is a
free module with a basis element of degree $r+1$. We saw in
\hyperref[fac:III:3:no:62]{no.~62} that $H^r(\Omega)=H_k^r(\Omega)$ for
all sufficiently large $k$, and that $H_k^r(\Omega)$ has over $K$ the
basis consisting of the single element
\phantomsection\label{fac:III:4:no:71:formula:xi-basis}
\[
\frac{(t_0\cdots t_r)^k}{t_0\cdots t_r}.
\]
The image of this element in $H^r(\Omega)$ will be denoted by $\xi$;
thus $\xi$ is a basis of $H^r(\Omega)$.

We shall now define a scalar product $\langle h,\varphi\rangle$ between
elements
\[
h\in B^r(M)_{-n}
\qquad\text{and}\qquad
\varphi\in\operatorname{Hom}_S(M,\Omega)_n,
\]
where $M$ is an arbitrary graded $S$-module. The element $\varphi$ can
be identified with an element of
$\operatorname{Hom}_S(M(-n),\Omega)_0$, that is, with a homomorphism
from $M(-n)$ to $\Omega$. Passing to cohomology groups, it therefore
defines a homomorphism
\[
H^r(M(-n))=B^r(M)_{-n}\longrightarrow H^r(\Omega),
\]
which we shall again denote by $\varphi$. The image of $h$ under this
homomorphism is a scalar multiple of $\xi$, and we define
$\langle h,\varphi\rangle$ by the formula
\phantomsection\label{fac:III:4:no:71:formula:pairing}
\[
\varphi(h)=\langle h,\varphi\rangle\,\xi.
\]

For every $\varphi\in\operatorname{Hom}_S(M,\Omega)_n$, the function
$h\mapsto\langle h,\varphi\rangle$ is a linear form on
$B^r(M)_{-n}$, and hence can be identified with an element
$\nu(\varphi)$ of the dual of $B^r(M)_{-n}$, which is precisely
$T^r(M)_n$. We have thus defined a homogeneous map of degree zero
\phantomsection\label{fac:III:4:no:71:map:nu}
\[
\nu:\operatorname{Hom}_S(M,\Omega)\longrightarrow T^r(M),
\]
and the formula
\phantomsection\label{fac:III:4:no:71:formula:S-linearity}
\[
\langle Ph,\varphi\rangle=\langle h,P\varphi\rangle
\]
shows that $\nu$ is an $S$-homomorphism.

\phantomsection\label{fac:III:4:no:71:prop:2}
\textsc{Proposition 2.} \textit{The homomorphism
$\nu:\operatorname{Hom}_S(M,\Omega)\longrightarrow T^r(M)$ is
bijective.}

\phantomsection\label{fac:III:4:no:71:proof:free-case}
We first prove \hyperref[fac:III:4:no:71:prop:2]{the proposition} when $M$
is a \textit{free} module. If
$M$ is a direct sum of homogeneous submodules $M^\alpha$, then
\phantomsection\label{fac:III:4:no:71:formula:direct-products}
\[
\operatorname{Hom}_S(M,\Omega)_n
=\prod_\alpha\operatorname{Hom}_S(M^\alpha,\Omega)_n,
\qquad
T^r(M)_n=\prod_\alpha T^r(M^\alpha)_n.
\]
Thus, if \hyperref[fac:III:4:no:71:prop:2]{the proposition} is true for the
$M^\alpha$, it is true for
$M$. This reduces the case of free modules to that of a free module on
one generator, that is, to the case $M=S(m)$.

We can then identify $\operatorname{Hom}_S(M,\Omega)_n$ with
$\operatorname{Hom}_S(S,S(n-m-r-1))_0$, that is, with the vector space
of homogeneous polynomials of degree $n-m-r-1$. Hence
$\operatorname{Hom}_S(M,\Omega)_n$ has as a basis the monomials
\phantomsection\label{fac:III:4:no:71:formula:polynomial-basis}
\[
t_0^{\gamma_0}\cdots t_r^{\gamma_r},
\qquad
\gamma_i\geq0,
\qquad
\sum_{i=0}^{i=r}\gamma_i=n-m-r-1.
\]
On the other hand, we saw in \hyperref[fac:III:3:no:62]{no.~62} that
$H_k^r(S(m-n))$ has as a
basis, for all sufficiently large $k$, the monomials
\[
\frac{(t_0\cdots t_r)^k}
     {t_0^{\beta_0}\cdots t_r^{\beta_r}},
\qquad
\beta_i>0,
\qquad
\sum_{i=0}^{i=r}\beta_i=n-m.
\]
Putting $\beta_i=\gamma_i'+1$, these monomials can be written as
\phantomsection\label{fac:III:4:no:71:formula:cohomology-basis}
\[
\frac{(t_0\cdots t_r)^{k-1}}
     {t_0^{\gamma_0'}\cdots t_r^{\gamma_r'}},
\qquad
\gamma_i'\geq0,
\qquad
\sum_{i=0}^{i=r}\gamma_i'=n-m-r-1.
\]
Returning to
\hyperref[fac:III:4:no:71:formula:pairing]{the definition of
$\langle h,\varphi\rangle$}, we find
that the scalar product
\phantomsection\label{fac:III:4:no:71:formula:dual-basis-pairing}
\[
\left\langle
 \frac{(t_0\cdots t_r)^{k-1}}
      {t_0^{\gamma_0'}\cdots t_r^{\gamma_r'}},
 t_0^{\gamma_0}\cdots t_r^{\gamma_r}
\right\rangle
\]
is always zero unless $\gamma_i=\gamma_i'$ for every $i$, in which
case it is equal to $1$. This means that $\nu$ sends the basis of the
$t_0^{\gamma_0}\cdots t_r^{\gamma_r}$ to the basis dual to the basis
of the
$(t_0\cdots t_r)^{k-1}/(t_0^{\gamma_0}\cdots t_r^{\gamma_r})$.
Consequently $\nu$ is bijective, which proves
\hyperref[fac:III:4:no:71:prop:2]{the proposition} when $M$ is free.

\phantomsection\label{fac:III:4:no:71:proof:general-case}
We now pass to the general case. Choose an exact sequence
\phantomsection\label{fac:III:4:no:71:formula:free-presentation}
\[
L^1\longrightarrow L^0\longrightarrow M\longrightarrow0
\]
in which $L^0$ and $L^1$ are free. Consider the following commutative
diagram:
\phantomsection\label{fac:III:4:no:71:diagram:nu}
\begin{center}
\begin{tikzcd}
0 \arrow[r] \arrow[d,"\nu"']
  & \operatorname{Hom}_S(M,\Omega) \arrow[r] \arrow[d,"\nu"']
  & \operatorname{Hom}_S(L^0,\Omega) \arrow[r] \arrow[d,"\nu"']
  & \operatorname{Hom}_S(L^1,\Omega) \arrow[d,"\nu"'] \\
0 \arrow[r]
  & T^r(M) \arrow[r]
  & T^r(L^0) \arrow[r]
  & T^r(L^1).
\end{tikzcd}
\end{center}

The first row of \hyperref[fac:III:4:no:71:diagram:nu]{this diagram} is exact
by \hyperref[fac:III:4:no:68:properties]{the general properties of the
functor $\operatorname{Hom}_S$}; the second row is also exact,
because it is dual to the sequence
\phantomsection\label{fac:III:4:no:71:formula:B-exact}
\[
B^r(L^1)\longrightarrow B^r(L^0)\longrightarrow B^r(M)
\longrightarrow0,
\]
which is itself exact by
\hyperref[fac:III:4:no:70:formula:long-exact-B]{the exact cohomology
sequence for the $B^q$} and
\hyperref[fac:III:4:no:70:property:B-vanishing]{the fact that
$B^{r+1}(M)=0$} for every graded $S$-module $M$.
Moreover, the two vertical homomorphisms
\phantomsection\label{fac:III:4:no:71:formula:free-nu}
\[
\nu:\operatorname{Hom}_S(L^0,\Omega)\longrightarrow T^r(L^0)
\quad\text{and}\quad
\nu:\operatorname{Hom}_S(L^1,\Omega)\longrightarrow T^r(L^1)
\]
are bijective,
\hyperref[fac:III:4:no:71:proof:free-case]{as we have just seen}. It follows
that
\phantomsection\label{fac:III:4:no:71:formula:general-nu}
\[
\nu:\operatorname{Hom}_S(M,\Omega)\longrightarrow T^r(M)
\]
is also bijective, which completes
\hyperref[fac:III:4:no:71:prop:2]{the proof}.
